\chapter{Introduction}
\label{intro}

Recent developments led to the emergence of the field of levitodynamics.
This field is primarily concerned with understanding, controlling, and using levitated nanoparticles.
It does so by using techniques developed in other research fields such as atomic physics and biophysics \cite{doi:10.1126/science.abg3027}.
A large contributor was 2018 Nobel Prize winner Arthur Ashkin \cite{10.1063/1.88748,PhysRevLett.24.156},
    whose early work on optical tweezers is resurfacing in this field \cite{PhysRevLett.122.123601}.
His research was continued 
    and current experimental work in this field has further advanced the frontier \cite{troyer2025quantumgroundstatecoolinglibrational}.
The high finesse cavities have also developed in recent times \cite{PhysRevLett.133.233603,Millen_2020}.
Levitodynamics has already been successful in the production of precise and novel sensors \cite{Moore_2021,doi:10.1126/sciadv.abo2361};
    in the future, a large goal is, among others, to gain insight into the quantum-to-classical transition \cite{gupta2025quantumtheoryelectricallylevitated,Millen_2020}.
\par 
Our model is based on References \cite{Maurer_2023,Gonzalez_Ballestero_2019}.
In this thesis, we investigate a trapped nanoparticle
    that is coupled to an optical cavity.
This cavity is driven by the scattering of photons off the particle.
A sketch of the setup can be seen in Figure \ref{fig:sp:setup}.
We then find the equilibrium positions by analytic and numerical methods.
For the numeric part, we are using the Julia Programming language \cite{Julia-2017},
    which provides high performance and high adaptability.
Contrary to the research we base this thesis on, 
    we treat this system classically.
Establishing a rigorous baseline is a necessary step in identifying quantum mechanical effects in the future.
Since we do not explicitly study related quantum phenomena,
    we may restrict ourselves to the classical description.
%We would expect that the quantum expected values will follow the classical results,
%    following from the Ehrenfest theorem.
\par
In Chapter \ref{ch:sp} we will derive the model for a single particle 
    and give the equilibrium positions as a plot of the detuning vs the $z$ values.
Using the numeric results we use the Hessian to compute the stability of the system.
\par
In Chapter \ref{ch:np} we extend this model to include $N$-particles. 
Here we also introduce necessary conditions 
    and generalize the model.
\par
In Chapter \ref{ch:2p} we use this model to find the equilibrium positions for $2$ particles.
We generate and analyze those results for different phases in the tweezers.
\par
In Chapter \ref{ch:sum} we summarize our results 
    and give a brief overview.
\par
Finally, in Chapter \ref{ch:otl} we give an outlook on future research options 
    and topics worth analyzing further.
